\newpage
\section{Khai báo về việc sử dụng AI trong dự án}

Trong quá trình thực hiện đồ án, nhóm có sử dụng công cụ Trí tuệ Nhân tạo (Gemini) như một nguồn hỗ trợ để tối ưu hóa quy trình làm việc và nâng cao chất lượng báo cáo. Việc sử dụng AI được khai báo minh bạch và cụ thể như sau:

\begin{itemize}
    \item \textbf{Mục đích sử dụng:} AI được dùng để gợi ý ý tưởng, tham khảo tài liệu, hỗ trợ viết nội dung báo cáo, và kiểm tra ngữ pháp/ngôn ngữ.
    \item \textbf{Phạm vi đóng góp:} 
    \begin{enumerate}
        \item Gợi ý cấu trúc báo cáo, bố cục các phần mục và cách trình bày chuyên nghiệp.
        \item Hỗ trợ tìm kiếm tài liệu tham khảo liên quan đến chủ đề quản trị nội dung (CMS) và bảo mật Web.
        \item Hỗ trợ kiểm tra ngữ pháp, chính tả và cải thiện câu văn trong các đoạn mô tả kỹ thuật.
        \item Đề xuất các công nghệ phù hợp với yêu cầu (Native PHP, MySQL, kiến trúc MVC).
        \item Hỗ trợ xây dựng định hướng thiết kế hệ thống và luồng xử lý dữ liệu qua Router.
        \item Phân tích yêu cầu chức năng (Task 1: Liên hệ, Task 2: Q\&A) và các yêu cầu phi chức năng.
        \item Viết nháp cho một số đoạn mô tả kỹ thuật về mô hình MVC, cơ chế PDO và xử lý bảo mật cơ bản.
        \item Gợi ý ý tưởng xây dựng cho các biểu đồ UML và mockup UI dựa trên mô tả của nhóm.
        \item Xây dựng một số đoạn mã mẫu minh họa cho logic xử lý Controller và Model.
    \end{enumerate}
    \item \textbf{Giới hạn:} AI không được sử dụng để tạo ra toàn bộ nội dung báo cáo và không thay thế cho tư duy phân tích của thành viên. Tất cả các thành phần cốt lõi như Use-case, lược đồ ERD, thiết kế Class và Mockup UI thực tế đều do nhóm tự xây dựng và chỉnh sửa.
    \item \textbf{Trách nhiệm:} Nhóm chịu trách nhiệm hoàn toàn về nội dung đã nộp, đảm bảo hiểu rõ và làm chủ toàn bộ mã nguồn cũng như các phần có sự hỗ trợ từ AI.
\end{itemize}

\noindent Việc sử dụng AI chỉ mang tính chất tham khảo và hỗ trợ, không thay thế cho năng lực phân tích, sáng tạo và đóng góp cá nhân của các thành viên trong nhóm.

\section{Mã nguồn dự án}

Để phục vụ công tác kiểm tra và đánh giá, nhóm đã lưu trữ toàn bộ tài liệu demo và mã nguồn dự án tại các liên kết dưới đây:

\begin{itemize}
    \item \textbf{Kho lưu trữ mã nguồn (GitHub):} \href{https://github.com/NghiaNguyenhk2005/WEB252}{Truy cập Repository dự án WEB252}
\end{itemize}

\textit{Lưu ý: Mã nguồn trên GitHub bao gồm đầy đủ file cấu hình \texttt{.htaccess}, file \texttt{seed.sql} để khởi tạo cơ sở dữ liệu và các thư mục theo kiến trúc MVC đã trình bày trong báo cáo.}